% 連続体力学 連続体力学-2 問題3（011-force） 解いた日 2026-10-24
%% kotae
\toi{1}{◎}{}
体：a，c，d

面：b，e
\toi{2}{◎}{}
表面をつらぬく力

内力どうしは作用・反作用で打ち消し合い，表面をつらぬく方向の力のみが残るため，表面積に比例する。
\toi{3}{△}{未着手}
面：$\ell^{2}$

体：$\ell^{3}$

\nuke{$\ell\to0$ の説明}
\naze{問いの後半が無い。積分形の運動方程式を $\int_{S}\vect{p}\,\dd S=\int_{V}(\rho\,\mathrm{D}\vect{u}/\mathrm{D}t-\vect{K})\,\dd V$ と移項し，右辺の被積分関数の大きさが $M$ 以下なら右辺の大きさは $M\ell^{3}$ 以下である。両辺を $\ell^{2}$ で割ると右辺は $M\ell$ 以下になり，$\ell\to0$ で $0$ に近づく。}
%% kizuki
\item 「$\ell$ の何乗か」は説明の材料で，それを使って移項した運動方程式 $\int_{S}\vect{p}\,\dd S=\int_{V}(\rho\,\mathrm{D}\vect{u}/\mathrm{D}t-\vect{K})\,\dd V$ の右辺を $M\ell^{3}$ で抑え，$\ell^{2}$ で割って $M\ell\to0$，まで書いて答えになる。
\item \mae{第1回 問題4 (4)・10-20} 小問の後半の「説明せよ」を書かずに止めた（前回は「理由」の部分が抜けた）。答える前に問題文の「求めよ」「説明せよ」を数え，答えの数とそろえる。
